\chapter{Referencial Teórico}
\label{chap:referencial}

Este capítulo apresenta o conjunto de conceitos fundamentais necessários para
compreender a análise conduzida nesta dissertação. São discutidos: (i) a
formação da volatilidade implícita em mercados de opções; (ii) a mensuração da
volatilidade realizada; (iii) a definição, interpretações e propriedades do
\textit{Volatility Risk Premium} (VRP); (iv) peculiaridades do mercado de opções
de Bitcoin; (v) regimes de volatilidade e sua relação com o BVRP; e
(vi) abordagens de aprendizado de máquina --- incluindo modelos lineares
regularizados (LASSO e Ridge) e métodos baseados em árvores (\textit{Random
Forest} e \textit{Gradient Boosting}) --- aplicadas à previsão do BVRP e de
retornos do Bitcoin.

\section{Volatilidade implícita}

A volatilidade implícita (\textit{Implied Volatility}, IV) representa a
expectativa de mercado sobre a variabilidade futura do preço do ativo
subjacente, inferida a partir dos preços das opções pela inversão de modelos de
precificação, como o modelo de Black--Scholes. Em vez de assumir uma volatilidade futura fixa, o mercado revela,
por meio dos preços de opções, a volatilidade que equilibra oferta e demanda por
proteção.

Diferentemente da volatilidade histórica, a volatilidade implícita não é uma
estimativa estatística no sentido clássico. Ela pode ser interpretada como um
\textit{preço} de risco associado à distribuição risco-neutra (\textit{risk-neutral})
dos retornos, refletindo não apenas a expectativa sobre a variância futura, mas
também prêmios por assimetria, curtose e risco de cauda. Assim, a IV é
influenciada por fatores como:

\begin{itemize}
    \item \textbf{demanda por proteção}: investidores adquirindo opções de venda
          (\textit{puts}) tendem a elevar a IV, especialmente em \textit{strikes}
          fora do dinheiro;
    \item \textbf{estrutura de \textit{skew}/\textit{smile}}: assimetrias de
          oferta e demanda entre diferentes \textit{strikes} e vencimentos;
    \item \textbf{eventos esperados}: anúncios regulatórios, atualizações de rede,
          vencimentos relevantes e mudanças de regime percebidas;
    \item \textbf{condições de liquidez}: \textit{bid--ask spreads}, profundidade
          do livro de ofertas e variações intradiárias de fluxo.
\end{itemize}

Em mercados tradicionais, a superfície de volatilidade apresenta padrões
estilizados bem documentados, como \textit{smile} e \textit{skew}, associados a
distribuições assimétricas e a risco de cauda. No caso do Bitcoin, a curva de
volatilidade implícita tende a apresentar inclinação mais acentuada e maior
sensibilidade a fluxos de curto prazo, refletindo a elevada incerteza do ativo,
a negociação contínua e a menor maturidade relativa do mercado, em linha com a
literatura sobre superfícies de volatilidade \citep{gatheral2011volsurface,
bakshi2003dynamicsIV}.

\section{Volatilidade realizada}

A volatilidade realizada (\textit{Realized Volatility}, RV) é uma medida
\textit{ex post} da variabilidade dos retornos observados de um ativo ao longo
do tempo, construída sob a medida histórica (física). Trata-se de um estimador
não paramétrico da volatilidade, amplamente utilizado na literatura empírica,
e constitui um insumo fundamental para a comparação com a volatilidade implícita.

Neste trabalho, a volatilidade realizada é calculada a partir de retornos
logarítmicos diários, considerando uma janela móvel de \(h\) dias, conforme:

\[
RV_t^{(h)} = \sqrt{ \frac{365}{h} \sum_{j=0}^{h-1} r_{t-j}^2 },
\]

onde \(r_{t-j} = \ln\left(\frac{P_{t-j}}{P_{t-j-1}}\right)\) representa o retorno
logarítmico diário do ativo, e o fator de anualização por 365 dias reflete a
negociação contínua do mercado de Bitcoin.

Alternativamente, pode-se definir a variância realizada como:

\[
\mathrm{VarRV}_t^{(h)} = \frac{365}{h} \sum_{j=0}^{h-1} r_{t-j}^2,
\]

sendo a volatilidade realizada a raiz quadrada dessa quantidade.

Embora a volatilidade realizada seja usualmente construída a partir de retornos
logarítmicos simples, pode-se alternativamente considerar retornos em excesso
em relação à taxa livre de risco, definidos como:

\[
r_t^e = r_t - r_{f,t},
\]

onde \(r_{f,t}\) representa a taxa livre de risco no período \(t\).

Nesse caso, a volatilidade realizada pode ser definida como:

\[
RV_t^{(h)} = \sqrt{ \frac{365}{h} \sum_{j=0}^{h-1} \left( r_{t-j} - r_{f,t-j} \right)^2 }.
\]

No entanto, para frequências diárias e, em particular, no contexto de criptoativos,
a diferença empírica entre utilizar retornos simples ou em excesso tende a ser
negligenciável. Por essa razão, a literatura usualmente adota retornos não ajustados,
abordagem também seguida neste trabalho.

Embora a RV possa ser estimada com dados intradiários, aumentando a precisão
(\textit{high-frequency realized volatility}), nesta dissertação utiliza-se a
versão baseada em dados diários, consistente com a disponibilidade de dados do
mercado à vista e com as séries derivadas do ecossistema da Deribit, incluindo
o índice DVOL.

A literatura clássica de volatilidade realizada, especialmente os trabalhos de
\citet{andersen2003realized} e \citet{corsi2009har}, estabelece a base teórica
para a construção de estimadores robustos de volatilidade, além de documentar a
presença de persistência temporal e estrutura de memória longa na dinâmica da
volatilidade, fenômenos também observados em criptoativos.
\section{Volatility Risk Premium (VRP)}

O \textit{Volatility Risk Premium} (VRP) pode ser interpretado como a
compensação exigida pelos investidores para assumir risco de volatilidade. Em
termos empíricos, é comum aproximar o VRP pela diferença entre uma medida de
volatilidade implícita e uma medida de volatilidade realizada no mesmo horizonte:

\[
VRP_t = IV_{30,t} - RV_{30,t},
\]

onde:

\begin{itemize}
    \item $IV_{30,t}$ é a volatilidade implícita a 30 dias (annualizada);
    \item $RV_{30,t}$ é a volatilidade realizada a 30 dias.
\end{itemize}

Um VRP positivo indica que agentes de mercado estão dispostos a pagar um prêmio
pela volatilidade implícita acima da volatilidade que efetivamente se realiza,
sendo usualmente interpretado como um \textit{prêmio de proteção}. Em mercados
tradicionais, como o S\&P 500, a evidência empírica sugere que o VRP apresenta
comportamento sistematicamente positivo, associado à demanda por proteção contra
quedas abruptas e contra choques de volatilidade \citep{bollerslev2011VRP,
pan2006jumpRisk, bali2009volatilitypremium}.

Do ponto de vista teórico, a distinção entre medida risco-neutra (associada a
preços de opções) e medida histórica (associada a retornos observados) implica
que a diferença IV--RV incorpora, além de expectativas de variância, componentes
relacionados a risco de cauda e a preferências por proteção. Dessa forma, o VRP
tem papel relevante tanto na precificação de ativos quanto em aplicações
empíricas de previsão.

No Bitcoin, entretanto, o VRP tende a ser mais instável. Entre os fatores
frequentemente associados a essa instabilidade, destacam-se:

\begin{enumerate}
    \item volatilidade realizada mais elevada, com episódios intermitentes de
          grande magnitude;
    \item eventos idiossincráticos e choques de liquidez, com impacto rápido na
          volatilidade;
    \item estrutura de mercado ainda em consolidação, com participação
          institucional relativamente menor e maior relevância de agentes
          alavancados.
\end{enumerate}

Esses elementos contribuem para maior dispersão do VRP e para ocorrência mais
frequente de períodos de VRP negativo, particularmente em janelas de realização
excepcionalmente elevada da volatilidade.

\section{Mercado de opções de Bitcoin}

O mercado de opções de Bitcoin é amplamente dominado pela Deribit, que oferece
liquidez superior às demais plataformas. Suas principais características incluem:

\begin{itemize}
    \item \textbf{negociação 24/7}: ausência de janelas de fechamento reduz o
          acúmulo de eventos fora do pregão, mas aumenta a relevância de efeitos
          intradiários de liquidez;
    \item \textbf{liquidez concentrada}: maior profundidade em opções
          \textit{at-the-money} e vencimentos curtos (7--30 dias), o que influencia
          a formação de IV em horizontes próximos;
    \item \textbf{efeito de liquidações automáticas}: posições alavancadas podem
          gerar movimentos bruscos e amplificar a volatilidade;
    \item \textbf{índice DVOL}: medida padronizada de volatilidade implícita
          calculada a partir de opções da Deribit.
\end{itemize}

A formação da IV no Bitcoin é sensível a fatores microestruturais, tais como:

\begin{itemize}
    \item tamanho mínimo de lote e restrições operacionais;
    \item \textit{bid--ask spreads} mais amplos em horários de menor profundidade;
    \item heterogeneidade de \textit{market makers} por região e perfil de risco;
    \item eventos associados ao \textit{funding} de futuros perpétuos, que afetam
          posicionamento e demanda por \textit{hedge}.
\end{itemize}

Essas características tornam o BVRP um objeto empírico mais rico, porém mais
volátil, do que em mercados tradicionais. A sensibilidade da volatilidade
implícita à microestrutura está alinhada com estudos sobre superfícies de
volatilidade e dinâmica de \textit{implied vol} \citep{gatheral2011volsurface,
bakshi2003dynamicsIV}, que destacam como liquidez, assimetria de fluxos e
microestrutura afetam a formação da curva de IV, especialmente em mercados menos
maduros ou mais fragmentados.

\section{Regimes de volatilidade}

Modelos de regimes buscam capturar mudanças estruturais no comportamento da
volatilidade e do VRP. Em mercados de criptoativos, essas transições tendem a ser
mais frequentes, refletindo a ocorrência de choques e mudanças rápidas no
sentimento de mercado. Regimes típicos incluem:

\begin{enumerate}
    \item \textbf{regime de baixa volatilidade}: períodos de lateralização do
          preço e menor demanda por proteção;
    \item \textbf{regime de alta volatilidade}: movimentos abruptos, liquidações
          em cascata e aumento expressivo do custo de \textit{hedge};
    \item \textbf{regime de transição}: aumento gradual da incerteza, com elevação
          progressiva de IV e RV.
\end{enumerate}

Mesmo sem a implementação explícita de modelos de mudança de regime, como
\textit{Markov Switching}, a inspeção gráfica e análises por quantis do BVRP
frequentemente revelam padrões compatíveis com tais regimes, motivando o uso de
estratégias condicionais e especificações empíricas que permitam relações
não lineares.

\section{Aprendizado de máquina em finanças}

Modelos de \textit{machine learning} têm se tornado ferramentas cada vez mais
comuns na pesquisa empírica em finanças, especialmente em aplicações de previsão
de retornos e volatilidade. Em contextos com múltiplas variáveis explicativas,
possíveis não linearidades e interações complexas, métodos flexíveis podem
superar especificações lineares tradicionais. Trabalhos como \citet{Gu2020EmpiricalML}
mostram que modelos baseados em árvores, redes neurais e técnicas de regularização
são capazes de capturar relações não lineares e de alta dimensão entre
características (\textit{features}) e retornos de ativos.

Nesta dissertação, quatro classes de modelos supervisionados são empregadas:
modelos lineares regularizados (LASSO e Ridge), que servem como referência
intermediária entre OLS e abordagens não lineares; e métodos baseados em árvores
(\textit{Random Forest} e \textit{Gradient Boosting}), utilizados pela capacidade
de capturar não linearidades, interações entre variáveis e de quantificar a
importância relativa de preditores como o BVRP em um ambiente multivariado.

\subsection{Modelos lineares regularizados: LASSO e Ridge}

Técnicas de regularização linear introduzem penalidades sobre os coeficientes
estimados com o objetivo de reduzir \textit{overfitting} e melhorar a
generalização fora da amostra, especialmente em contextos com múltiplos
preditores potencialmente correlacionados \citep{tibshirani1996lasso}.

O LASSO (\textit{Least Absolute Shrinkage and Selection Operator}) minimiza a
soma dos quadrados dos resíduos sujeita a uma penalidade $\ell_1$:
\[
\hat{\beta}^{\text{LASSO}} = \arg\min_{\beta}
    \left\{ \sum_{t=1}^{T}(y_t - \mathbf{x}_t'\beta)^2
            + \lambda \sum_{j=1}^{p} |\beta_j| \right\},
\]
onde $\lambda \geq 0$ é o parâmetro de regularização. A penalidade $\ell_1$
induz esparsidade: coeficientes associados a preditores irrelevantes são levados
exatamente a zero, realizando seleção automática de variáveis.

O Ridge utiliza penalidade $\ell_2$:
\[
\hat{\beta}^{\text{Ridge}} = \arg\min_{\beta}
    \left\{ \sum_{t=1}^{T}(y_t - \mathbf{x}_t'\beta)^2
            + \lambda \sum_{j=1}^{p} \beta_j^2 \right\}.
\]
Diferentemente do LASSO, o Ridge retém todos os preditores, mas encolhe seus
coeficientes em direção a zero, sendo especialmente útil na presença de
multicolinearidade. Nesta dissertação, ambos os modelos funcionam como
\textit{benchmarks} lineares robustos, intermediários entre a regressão OLS
irrestrita e as abordagens baseadas em árvores.

\subsection{Random Forest}

O Random Forest é um método de \textit{bagging} que cria múltiplas árvores de
decisão \citep{breiman2001randomforest}, cada uma treinada em uma amostra
\textit{bootstrap} dos dados. Trata-se de um dos algoritmos mais consagrados da
literatura de aprendizado de máquina, amplamente utilizado em aplicações
financeiras devido à sua capacidade de capturar relações não lineares e
interações complexas entre variáveis.

Árvores agregadas em estrutura de \textit{bagging} apresentam vantagens como:

\begin{itemize}
    \item captura de não linearidades;
    \item robustez a outliers e ruído;
    \item bom desempenho mesmo quando relações lineares são fracas.
\end{itemize}

Além disso, a aleatorização de subconjuntos de variáveis em cada divisão reduz a
correlação entre árvores e tende a melhorar a generalização fora da amostra.

\subsection{Gradient Boosting}

\textit{Gradient Boosting} é uma técnica de \textit{ensemble} que constrói
árvores de decisão sequencialmente, de modo que cada nova árvore busca corrigir
os erros residuais das árvores anteriores \citep{friedman2001gradientboosting}.
Diferentemente do Random Forest, que utiliza \textit{bagging} com árvores
independentes, o \textit{Gradient Boosting} realiza otimização sequencial sobre
uma função de perda, resultando em maior sensibilidade a padrões sutis e
relações não lineares.

Nesta dissertação, utiliza-se a implementação XGBoost (\textit{eXtreme Gradient
Boosting}), introduzida por \citet{chen2016xgboost}, que combina alto desempenho
preditivo com regularização explícita e eficiência computacional. Entre suas
características distintivas destacam-se:

\begin{itemize}
    \item construção sequencial de árvores para correção iterativa dos resíduos;
    \item regularização L1 e L2 integrada, controlando \textit{overfitting};
    \item ajuste fino do viés e da variância por meio de hiperparâmetros como
          profundidade máxima e taxa de aprendizado (\textit{learning rate});
    \item implementação eficiente e paralelizável, adequada a séries temporais
          de dimensionalidade moderada.
\end{itemize}

Os quatro modelos --- LASSO, Ridge, \textit{Random Forest} e \textit{Gradient
Boosting} --- são empregados nesta dissertação para avaliar se o BVRP e variáveis
de estado de mercado contêm informação preditiva relevante sobre o próprio prêmio
e sobre retornos futuros do Bitcoin, sob protocolos rigorosos de validação fora
da amostra.

\section{Síntese do capítulo}

Este capítulo apresentou os conceitos fundamentais que sustentam a investigação
sobre o prêmio de risco de volatilidade no Bitcoin. Foram revisados elementos
teóricos da volatilidade implícita e realizada, o cálculo e a interpretação do
VRP, aspectos microestruturais específicos do mercado de opções de Bitcoin, a
identificação de regimes de volatilidade e a motivação para o uso de quatro
classes de modelos de aprendizado de máquina (LASSO, Ridge, \textit{Random
Forest} e \textit{Gradient Boosting}) na previsão do BVRP e de retornos futuros.

Esses elementos formam a base conceitual para os dados e a metodologia
apresentados nos Capítulos~\ref{chap:dados} e~\ref{chap:metodologia},
nos quais são detalhadas as fontes de dados, o processo de construção das
variáveis de interesse e a especificação dos modelos empíricos utilizados
ao longo da dissertação.
